\chapter{Database Schema Design}
\chaptermark{Database Schema Design}
\label{ch:db-schema}

This chapter presents the database schema for the secure online banking platform. The schema is implemented in \textbf{PostgreSQL} and enforces data integrity, financial accuracy, and auditability at the storage layer. Seven interrelated tables capture every aspect of the system---from user identity and account management to transaction processing and compliance logging.

\section{Deployment platform}

The original project design considered hosting the relational database on \textbf{Microsoft Azure} (for example Azure SQL Database). During setup we were \textbf{unable to create or activate an Azure account} (registration, verification, or subscription constraints). To continue the coursework without blocking on that provider, we deployed the same logical schema on \textbf{PostgreSQL} using \textbf{Neon's free tier}. Neon provides managed PostgreSQL with TLS, a web SQL editor, and a connection string suitable for local development and demonstration. The DDL in \texttt{banking\_schema.sql} is written for PostgreSQL and matches the tables created on Neon.

\section{Design principles}

The schema was designed around four guiding principles that align with the broader banking security design (identity controls, RBAC, encryption in transit, and auditability):

\begin{description}
  \item[Referential integrity:]
    Foreign-key constraints link every account to a user, every transaction to its accounts, and every audit entry to the actor who performed it. Orphaned records are structurally impossible.
  \item[Double-entry bookkeeping:]
    Every financial movement produces two entries---a debit and a credit---ensuring that the ledger always balances and discrepancies are immediately detectable.
  \item[Least privilege and role separation:]
    The \texttt{users} table stores an RBAC role (\texttt{customer}, \texttt{teller}, \texttt{manager}, \texttt{admin}) that the backend API should enforce at every endpoint.
  \item[Immutable audit trail:]
    The \texttt{audit\_log} table captures before/after snapshots of state changes, supporting non-repudiation and compliance review.
\end{description}

\section{Entity--relationship diagram}

Figure~\ref{fig:er-diagram} shows the entity--relationship diagram exported from the schema design tool (dbdiagram.io), including tables and relationships as implemented in PostgreSQL.

\begin{figure}[htbp]
  \centering
  \includegraphics[width=\textwidth,height=0.88\textheight,keepaspectratio,page=1]{db_schema_diagram.pdf}
  \caption{Entity--relationship diagram of the banking schema (PostgreSQL).}
  \label{fig:er-diagram}
\end{figure}

\section{Relationship summary}

Table~\ref{tab:er-relationships} summarizes the relationships between the seven tables. Primary and foreign keys follow the naming used in \texttt{banking\_schema.sql}.

\begin{table}[htbp]
  \centering
  \small
  \setlength{\tabcolsep}{5pt}
  \renewcommand{\arraystretch}{1.15}
  \begin{tabularx}{\textwidth}{@{} l l X @{}}
    \toprule
    \textbf{Parent table} & \textbf{Child table} & \textbf{Relationship} \\
    \midrule
    users              & accounts            & One user owns zero or more accounts. \\
    users              & transactions        & One user initiates zero or more transactions. \\
    users              & audit\_log          & One user generates zero or more log entries. \\
    accounts           & transactions        & An account may be the source or destination of a transaction. \\
    accounts           & transaction\_entries & An account has zero or more ledger entries. \\
    accounts           & account\_limits     & An account has zero or more spending limits. \\
    transactions       & transaction\_entries & One transaction produces one or more ledger entries. \\
    transaction\_states & transactions       & One state applies to zero or more transactions. \\
    \bottomrule
  \end{tabularx}
  \caption{Relationship summary for the banking schema.}
  \label{tab:er-relationships}
\end{table}

\section{Table definitions}

Each subsection below describes one table: its purpose, columns, data types, constraints, and security rationale. The full SQL script is provided in the accompanying file \texttt{banking\_schema.sql}.

% ------------------------------------------------------------------
\subsection{users}
\label{sec:tbl-users}

The \texttt{users} table stores identity, authentication, and authorization data for every person who interacts with the platform.

\begin{table}[htbp]
  \centering
  \caption{\texttt{users} table columns (PostgreSQL)}
  \label{tab:users}
  \small
  \setlength{\tabcolsep}{4pt}
  \renewcommand{\arraystretch}{1.15}
  \begin{tabularx}{\textwidth}{@{} l l l X @{}}
    \toprule
    \textbf{Column} & \textbf{Type} & \textbf{Constraint} & \textbf{Purpose} \\
    \midrule
    user\_id            & UUID & PK, default \texttt{gen\_random\_uuid()} & Non-sequential surrogate key. \\
    first\_name         & VARCHAR(100) & NOT NULL     & Legal first name. \\
    last\_name          & VARCHAR(100) & NOT NULL     & Legal last name. \\
    email               & VARCHAR(255) & UNIQUE, NOT NULL & Login identifier. \\
    phone\_number       & VARCHAR(20)  & NULL         & Optional contact number. \\
    password\_hash      & BYTEA        & NOT NULL     & Bcrypt/Argon2 hash (never plaintext). \\
    password\_salt      & BYTEA        & NOT NULL     & Per-user cryptographic salt. \\
    date\_of\_birth     & DATE         & NOT NULL     & KYC age verification. \\
    national\_id        & VARCHAR(50)  & UNIQUE, NOT NULL & Government-issued ID for KYC. \\
    address fields      & VARCHAR      & NULL         & Line~1, Line~2, city, country, postal code. \\
    role                & VARCHAR(20)  & CHECK        & RBAC role: customer, teller, manager, admin. \\
    mfa\_enabled        & BOOLEAN      & DEFAULT FALSE & Whether MFA is active. \\
    mfa\_secret         & BYTEA        & NULL         & Encrypted TOTP seed. \\
    is\_active          & BOOLEAN      & DEFAULT TRUE & Soft-delete flag. \\
    is\_locked          & BOOLEAN      & DEFAULT FALSE & Set after repeated failed logins. \\
    failed\_login\_count& INT          & DEFAULT 0    & Counter for account-lockout policy. \\
    last\_login\_at     & TIMESTAMPTZ  & NULL         & Timestamp of most recent authentication. \\
    created\_at         & TIMESTAMPTZ  & DEFAULT \texttt{NOW()} & Row creation time. \\
    updated\_at         & TIMESTAMPTZ  & DEFAULT \texttt{NOW()} & Last modification time. \\
    \bottomrule
  \end{tabularx}
\end{table}

\noindent\textbf{Security notes.}
Passwords are stored as salted hashes using bcrypt or Argon2. The \texttt{is\_locked} flag and \texttt{failed\_login\_count} enforce account lockout, mitigating brute-force attacks. UUIDs are used instead of sequential integers to reduce predictable identifier enumeration.

% ------------------------------------------------------------------
\subsection{accounts}
\label{sec:tbl-accounts}

Each row in \texttt{accounts} represents a financial account owned by a user.

\begin{table}[htbp]
  \centering
  \caption{\texttt{accounts} table columns (PostgreSQL)}
  \label{tab:accounts}
  \small
  \setlength{\tabcolsep}{4pt}
  \renewcommand{\arraystretch}{1.15}
  \begin{tabularx}{\textwidth}{@{} l l l X @{}}
    \toprule
    \textbf{Column} & \textbf{Type} & \textbf{Constraint} & \textbf{Purpose} \\
    \midrule
    account\_id       & UUID & PK, default \texttt{gen\_random\_uuid()} & Non-sequential surrogate key. \\
    user\_id          & UUID & FK $\to$ \texttt{users} & Account owner. \\
    account\_number   & VARCHAR(20)     & UNIQUE        & Customer-facing identifier. \\
    account\_type     & VARCHAR(20)     & CHECK         & checking, savings, fixed\_deposit, loan. \\
    currency          & VARCHAR(3)      & DEFAULT USD   & ISO~4217 currency code. \\
    balance           & NUMERIC(18,2)   & DEFAULT 0     & Ledger (book) balance. \\
    available\_balance& NUMERIC(18,2)   & DEFAULT 0     & Balance minus pending holds. \\
    interest\_rate    & NUMERIC(5,4)    & NULL          & Annual rate for savings/loans. \\
    status            & VARCHAR(20)     & CHECK         & active, frozen, closed, dormant. \\
    opened\_at        & TIMESTAMPTZ     & DEFAULT \texttt{NOW()} & Account opening time. \\
    closed\_at        & TIMESTAMPTZ     & NULL          & Populated when account is closed. \\
    created\_at / updated\_at & TIMESTAMPTZ & DEFAULT \texttt{NOW()} & Row lifecycle timestamps. \\
    \bottomrule
  \end{tabularx}
\end{table}

\noindent\textbf{Security notes.}
The \texttt{status} column allows the system to freeze an account immediately if fraud is detected, without deleting data. Splitting \texttt{balance} and \texttt{available\_balance} prevents users from spending funds that are locked in pending transactions.

% ------------------------------------------------------------------
\subsection{transaction\_states}
\label{sec:tbl-txstates}

\texttt{transaction\_states} is a reference table that defines the finite set of lifecycle states a transaction can occupy.

\begin{table}[htbp]
  \centering
  \caption{\texttt{transaction\_states} table columns and seed data}
  \label{tab:txstates}
  \small
  \setlength{\tabcolsep}{4pt}
  \renewcommand{\arraystretch}{1.15}
  \begin{tabularx}{\textwidth}{@{} l l l X @{}}
    \toprule
    \textbf{Column} & \textbf{Type} & \textbf{Constraint} & \textbf{Purpose} \\
    \midrule
    state\_id    & INT          & PK, \texttt{GENERATED ALWAYS AS IDENTITY} & Surrogate key. \\
    state\_name  & VARCHAR(30)  & UNIQUE       & Machine-readable state label. \\
    description  & VARCHAR(255) & NULL         & Human-readable explanation. \\
    \bottomrule
  \end{tabularx}
\end{table}

\noindent The table is pre-seeded with seven states:

\begin{enumerate}
  \item \textbf{pending} --- transaction initiated, awaiting processing.
  \item \textbf{authorized} --- authorized but not yet settled.
  \item \textbf{completed} --- fully processed and settled.
  \item \textbf{failed} --- processing error occurred.
  \item \textbf{reversed} --- reversed after completion (e.g.\ chargeback).
  \item \textbf{cancelled} --- cancelled before settlement.
  \item \textbf{on\_hold} --- flagged for manual review by compliance.
\end{enumerate}

% ------------------------------------------------------------------
\subsection{transactions}
\label{sec:tbl-transactions}

The \texttt{transactions} table records every financial operation in the system.

\begin{table}[htbp]
  \centering
  \caption{\texttt{transactions} table columns (PostgreSQL)}
  \label{tab:transactions}
  \small
  \setlength{\tabcolsep}{4pt}
  \renewcommand{\arraystretch}{1.15}
  \begin{tabularx}{\textwidth}{@{} l l l X @{}}
    \toprule
    \textbf{Column} & \textbf{Type} & \textbf{Constraint} & \textbf{Purpose} \\
    \midrule
    transaction\_id    & UUID & PK, default \texttt{gen\_random\_uuid()} & Non-sequential surrogate key. \\
    reference\_number  & VARCHAR(50)     & UNIQUE           & External reference for receipts. \\
    transaction\_type  & VARCHAR(30)     & CHECK            & deposit, withdrawal, transfer, payment, fee, interest, refund. \\
    amount             & NUMERIC(18,2)   & CHECK $>$ 0      & Transaction value. \\
    currency           & VARCHAR(3)      & DEFAULT USD      & ISO~4217 code. \\
    description        & VARCHAR(500)    & NULL             & Free-text memo. \\
    state\_id          & INT             & FK $\to$ \texttt{transaction\_states} & Current lifecycle state. \\
    source\_account\_id& UUID            & FK $\to$ \texttt{accounts} & Debited account (NULL for deposits). \\
    dest\_account\_id  & UUID            & FK $\to$ \texttt{accounts} & Credited account (NULL for withdrawals). \\
    initiated\_by      & UUID            & FK $\to$ \texttt{users}   & Actor who started the transaction. \\
    ip\_address        & INET            & NULL             & Client IP (native PostgreSQL type). \\
    channel            & VARCHAR(20)     & CHECK            & web, mobile, atm, branch, api. \\
    scheduled\_at      & TIMESTAMPTZ     & NULL             & For future-dated transfers. \\
    processed\_at      & TIMESTAMPTZ     & NULL             & When settlement completed. \\
    created\_at / updated\_at & TIMESTAMPTZ & DEFAULT \texttt{NOW()} & Row lifecycle timestamps. \\
    \bottomrule
  \end{tabularx}
\end{table}

\noindent\textbf{Security notes.}
Capturing \texttt{ip\_address} and \texttt{channel} provides forensic context for fraud investigations. The \texttt{state\_id} foreign key ensures that only valid, pre-defined states can be assigned.

% ------------------------------------------------------------------
\subsection{transaction\_entries}
\label{sec:tbl-txentries}

\texttt{transaction\_entries} implements double-entry bookkeeping: every transaction generates at least two entries (a debit and a credit) so that the sum across all entries is always zero.

\begin{table}[htbp]
  \centering
  \caption{\texttt{transaction\_entries} table columns (PostgreSQL)}
  \label{tab:txentries}
  \small
  \setlength{\tabcolsep}{4pt}
  \renewcommand{\arraystretch}{1.15}
  \begin{tabularx}{\textwidth}{@{} l l l X @{}}
    \toprule
    \textbf{Column} & \textbf{Type} & \textbf{Constraint} & \textbf{Purpose} \\
    \midrule
    entry\_id        & BIGINT          & PK, \texttt{GENERATED ALWAYS AS IDENTITY} & Surrogate key. \\
    transaction\_id  & UUID            & FK $\to$ \texttt{transactions}  & Parent transaction. \\
    account\_id      & UUID            & FK $\to$ \texttt{accounts}      & Affected account. \\
    entry\_type      & VARCHAR(6)      & CHECK                  & \texttt{debit} or \texttt{credit}. \\
    amount           & NUMERIC(18,2)   & CHECK $>$ 0            & Entry value. \\
    balance\_before  & NUMERIC(18,2)   & NOT NULL               & Account balance before this entry. \\
    balance\_after   & NUMERIC(18,2)   & NOT NULL               & Account balance after this entry. \\
    created\_at      & TIMESTAMPTZ     & DEFAULT \texttt{NOW()} & Immutable creation timestamp. \\
    \bottomrule
  \end{tabularx}
\end{table}

\noindent\textbf{Security notes.}
Recording \texttt{balance\_before} and \texttt{balance\_after} creates an immutable ledger trail. If the running balance ever diverges from the sum of entries, the system can detect tampering or software bugs immediately.

% ------------------------------------------------------------------
\subsection{account\_limits}
\label{sec:tbl-limits}

\texttt{account\_limits} enforces configurable spending and transfer caps per account, providing a technical control against fraud and unauthorized fund movement.

\begin{table}[htbp]
  \centering
  \caption{\texttt{account\_limits} table columns (PostgreSQL)}
  \label{tab:limits}
  \small
  \setlength{\tabcolsep}{4pt}
  \renewcommand{\arraystretch}{1.15}
  \begin{tabularx}{\textwidth}{@{} l l l X @{}}
    \toprule
    \textbf{Column} & \textbf{Type} & \textbf{Constraint} & \textbf{Purpose} \\
    \midrule
    limit\_id      & INT             & PK, \texttt{GENERATED ALWAYS AS IDENTITY} & Surrogate key. \\
    account\_id    & UUID            & FK $\to$ \texttt{accounts}     & Account to which the limit applies. \\
    limit\_type    & VARCHAR(30)     & CHECK                 & daily\_withdrawal, daily\_transfer, single\_transaction, monthly\_withdrawal, monthly\_transfer. \\
    max\_amount    & NUMERIC(18,2)   & CHECK $>$ 0           & Maximum permitted value. \\
    current\_usage & NUMERIC(18,2)   & DEFAULT 0             & Running total for the current period. \\
    usage\_reset\_at& TIMESTAMPTZ    & NOT NULL              & When the counter resets. \\
    is\_active     & BOOLEAN         & DEFAULT TRUE          & Enable or disable the limit. \\
    created\_at / updated\_at & TIMESTAMPTZ & DEFAULT \texttt{NOW()} & Row lifecycle timestamps. \\
    \bottomrule
  \end{tabularx}
\end{table}

\noindent A composite unique constraint on \texttt{(account\_id, limit\_type)} ensures that each account can have at most one active limit of each kind.

% ------------------------------------------------------------------
\subsection{audit\_log}
\label{sec:tbl-audit}

The \texttt{audit\_log} table provides an immutable compliance trail for security events and data changes.

\begin{table}[htbp]
  \centering
  \caption{\texttt{audit\_log} table columns (PostgreSQL)}
  \label{tab:audit}
  \small
  \setlength{\tabcolsep}{4pt}
  \renewcommand{\arraystretch}{1.15}
  \begin{tabularx}{\textwidth}{@{} l l l X @{}}
    \toprule
    \textbf{Column} & \textbf{Type} & \textbf{Constraint} & \textbf{Purpose} \\
    \midrule
    log\_id          & BIGINT          & PK, \texttt{GENERATED ALWAYS AS IDENTITY} & Surrogate key. \\
    event\_type      & VARCHAR(50)     & NOT NULL        & Category (e.g.\ \texttt{user.login}). \\
    entity\_type     & VARCHAR(50)     & NOT NULL        & Table or resource affected. \\
    entity\_id       & VARCHAR(100)    & NOT NULL        & Primary key of the affected row. \\
    action           & VARCHAR(20)     & CHECK           & create, read, update, delete, login, logout, failed\_login. \\
    performed\_by    & UUID            & FK $\to$ \texttt{users}  & Actor (NULL for system events). \\
    old\_values      & JSONB           & NULL            & Snapshot before the change (queryable). \\
    new\_values      & JSONB           & NULL            & Snapshot after the change. \\
    ip\_address      & INET            & NULL            & Client IP for forensic correlation. \\
    user\_agent      & VARCHAR(500)    & NULL            & Browser or API client identifier. \\
    additional\_info & JSONB           & NULL            & Free-form context (e.g.\ error payloads). \\
    created\_at      & TIMESTAMPTZ     & DEFAULT \texttt{NOW()} & Immutable event timestamp. \\
    \bottomrule
  \end{tabularx}
\end{table}

\noindent\textbf{Security notes.}
Storing \texttt{old\_values} and \texttt{new\_values} as \texttt{JSONB} enables indexed queries and point-in-time reconstruction. In production, append-only storage and integrity hashing of log streams should complement this table design.

\section{Indexing strategy}

Performance-critical queries are supported by targeted indexes:

\begin{itemize}
  \item \textbf{accounts}: indexed on \texttt{user\_id} for fast lookup of all accounts belonging to a customer.
  \item \textbf{transactions}: indexed on \texttt{source\_account\_id}, \texttt{dest\_account\_id}, \texttt{state\_id}, and \texttt{created\_at} to accelerate statement generation, state-based filtering, and date-range queries.
  \item \textbf{transaction\_entries}: indexed on \texttt{transaction\_id} and \texttt{account\_id} for efficient ledger reconstruction.
  \item \textbf{account\_limits}: indexed on \texttt{account\_id} so that limit checks during transaction processing are fast.
  \item \textbf{audit\_log}: indexed on \texttt{(entity\_type, entity\_id)}, \texttt{performed\_by}, \texttt{created\_at}, and \texttt{action} to support compliance queries and incident investigation.
\end{itemize}

\section{Security alignment summary}

Table~\ref{tab:schema-security} maps schema features to the threats they help mitigate.

\begin{table}[htbp]
  \centering
  \caption{Schema features mapped to threats and controls}
  \label{tab:schema-security}
  \small
  \setlength{\tabcolsep}{5pt}
  \renewcommand{\arraystretch}{1.2}
  \begin{tabularx}{\textwidth}{@{} >{\raggedright\arraybackslash}p{3.8cm} >{\raggedright\arraybackslash}p{3.5cm} X @{}}
    \toprule
    \textbf{Schema feature} & \textbf{Threat addressed} & \textbf{Rationale} \\
    \midrule
    Salted password hashing       & Spoofing (credential theft)     & Passwords never stored in plaintext; slow hash resists offline guessing. \\
    MFA columns                   & Spoofing (session hijacking)    & Second factor reduces impact of stolen passwords. \\
    RBAC role column              & Elevation of privilege           & Enforced in application layer against role escalation. \\
    Account lockout fields        & Spoofing (brute force)          & Failed login counters and lock flags throttle guessing. \\
    Double-entry bookkeeping      & Tampering (balance manipulation)& Every movement has matching debit/credit entries. \\
    Transaction state machine     & Repudiation (state falsification)& States constrained by FK to reference data. \\
    Audit log with JSON snapshots & Repudiation (action denial)     & Actor, action, and before/after values recorded. \\
    IP and user-agent logging     & Forensics / incident response   & Correlates sessions with security events. \\
    Account limits                & Fraud / fund drain              & Caps velocity of withdrawals and transfers. \\
    UUID primary keys             & Enumeration                     & Non-sequential identifiers harder to guess than integers. \\
    \bottomrule
  \end{tabularx}
\end{table}
